\documentclass[10pt,a4paper]{article}

\usepackage[a4paper,left=0.35in,right=0.35in,top=0.27in,bottom=0.25in]{geometry}
\usepackage{fontspec}
\setmainfont{Arial}
\usepackage[english]{babel}
\usepackage{enumitem}
\usepackage{xcolor}
\usepackage{hyperref}
\usepackage{titlesec}
\usepackage{microtype}

\definecolor{primary}{HTML}{2463A5}
\definecolor{textdark}{HTML}{303030}
\definecolor{textlight}{HTML}{666666}
\definecolor{rulecolor}{HTML}{C7C7C7}
\hypersetup{colorlinks=true,urlcolor=primary,linkcolor=primary,
  pdftitle={Galaba Vamsi - Resume (Machine Learning)},pdfauthor={Galaba Vamsi}}

\pagestyle{empty}
\setlength{\parindent}{0pt}
\setlength{\parskip}{0pt}
\urlstyle{same}

% Body: 9.5 pt Arial on 11.4 pt leading, scaled by \linespread.
\linespread{0.90}
\renewcommand{\normalsize}{\fontsize{9.5}{11.4}\selectfont}

% No hyphenation: split tool names (MAT-LAB, Web-Sockets) break ATS keyword matching.
\hyphenpenalty=10000
\exhyphenpenalty=10000
\emergencystretch=1.5em
\frenchspacing

\titleformat{\section}
  {\fontsize{10.5}{12}\selectfont\bfseries\color{primary}}
  {}{0pt}{\MakeUppercase}
  [\vspace{1pt}{\color{rulecolor}\titlerule[0.45pt]}]
\titlespacing*{\section}{0pt}{3.4pt}{2.6pt}

\setlist[itemize]{
  label=\textbullet,
  leftmargin=1.0em,
  itemsep=0.08em,
  topsep=0.06em,
  parsep=0pt,
  partopsep=0pt
}

% \entry{bold title}{context and links}{dates}
\newcommand{\entry}[3]{%
  \par\vspace{0.16em}%
  \noindent{\color{textdark}\textbf{#1}}{\color{primary}\enspace#2}\hfill{\color{textlight}#3}\par
}

% \project{title}{context}{description with links}: a run-in title keeps each project short.
\newcommand{\project}[3]{%
  \par\vspace{0.1em}%
  \noindent\textbf{#1}{\color{textlight}\enspace(#2)}\enspace#3\par
}

% \edu{text}{dates}: dates flush right on the first line; the text wraps beneath them.
\newlength{\datewidth}
\newcommand{\edu}[2]{%
  \par\vspace{0.16em}%
  \settowidth{\datewidth}{\quad#2}%
  \parshape 2 0pt \dimexpr\linewidth-\datewidth\relax 0pt \linewidth
  \noindent\makebox[0pt][l]{\hspace{\dimexpr\linewidth-\datewidth\relax}\quad{\color{textlight}#2}}#1\par
}

\newcommand{\sep}{\enspace|\enspace}

% Arial maps ";" and U+037E to one glyph, so the PDF text layer reported Greek question
% marks. ActualText keeps copy/paste and ATS parsing faithful to the source characters.
\XeTeXgenerateactualtext=1

\begin{document}
\normalsize
\color{textdark}

% Single column, read top to bottom, so ATS parsers keep the order.
{\fontsize{21}{23}\selectfont\textbf{Galaba Vamsi}}\par
\vspace{0.25em}
{\color{primary}Mechatronics at IIT Bhilai. Machine learning, computer vision, and applied AI.}\par
{\color{textlight}
\href{tel:+917483047846}{+91 7483047846}\sep
\href{mailto:galabavamsi12@gmail.com}{galabavamsi12@gmail.com}\sep
\href{https://galabavamsi.github.io/portfolio}{galabavamsi.github.io/portfolio}\\
\href{https://www.linkedin.com/in/galabavamsi/}{linkedin.com/in/galabavamsi}\sep
\href{https://github.com/Galabavamsi}{github.com/Galabavamsi}}\par

\section*{Summary}
Graduating May 2027. I build and evaluate ML systems: wearables recognition (Kaggle, 5th of 232 teams), a voice-clone detector (hackathon finalist), vision-language price prediction, and I/Q signal classification. 15 merged pull requests to open-source robot-data and evaluation tools; I wrote and maintain inspect-robots-wandb on PyPI. Co-author of 3 papers.\par

\section*{Education}
\edu{\textbf{B.Tech. in Mechatronics},{\color{primary}\enspace Indian Institute of Technology Bhilai}\sep\textbf{Coursework:} AI and ML, Deep Learning for Computer Vision, Digital Image Processing, Optimal Control, Advanced Control Theory, DSA, OOP.}{Aug 2023 -- May 2027}

\section*{Projects}
\project{WEAR Dataset Challenge @HASCA 2026}{Kaggle}{Wearable exercise recognition from one-second IMU and video tiles of four unseen subjects (macro-F1 over 19 classes): 5th of 232 teams on the final private leaderboard (0.93137; 7th public, 0.93825). A fork of a public notebook; my additions took my reproduction from 0.92497 to 0.93825 public: per-subject whitening, whitened video-continuity link features, ZCA video input, and a multi-run ensemble decode. A global whitening that raised out-of-fold score but lowered the leaderboard score was a leak, and I replaced it with per-subject whitening. Built with Claude. \href{https://www.kaggle.com/competitions/3rd-wear-dataset-challenge-hasca-2026/leaderboard}{Leaderboard}}

\project{Pehchaan}{team of 2; top-30 finalist of 4,723, RAKSHAM (IIT Delhi $\times$ Amazon), 2026}{UPI voice-clone fraud defence; I led ML and backend.}
\begin{itemize}
  \item Fine-tuned XLS-R-300M + AASIST on consented Telugu, Hindi, and English speech and clones from 13 generators: 4.8\% EER vs 34.3\% off the shelf (6.0\% on 6 unseen generators).
  \item A 48-call real-carrier test exposed a gap (24.6\% EER). Learning the speakerphone channel and retraining cut it to 6.8\%, clones caught 48\% to 75\%, 0\% of genuine callers flagged. In simulation the risk engine holds 73\% of scams vs 0.21\% of genuine payments.
\end{itemize}

\project{Amazon ML Challenge}{team; 2025}{Product price prediction with a Qwen2-VL 7B and ViT/RoBERTa two-tower model; team rank 700 (20,000+ registered participants). Also AIR 200 in 2026.}

\project{Arista Wi-Fi RRM}{team; Inter IIT Tech Meet 14.0, 5th place, Dec 2025}{I/Q-based interference classification (STFT, DBSCAN, Dual-CNN; 16 classes) for client-aware radio resource management in NS-3 on a 5-AP, 50-client 802.11k/v topology. \href{https://canva.link/puoqbe4okoid8rb}{Slides}}

\project{Lunar Crater Detection}{2026}{Reproduced Zhou et al. (2018) on lunar DEMs: F1 0.839 on one labelled region (parameters tuned on that region). \href{https://github.com/Galabavamsi/Moon-Crater-Detection-using-DEM}{GitHub}}

\project{AntennaNet}{2026}{Inverse antenna design: surrogate network plus KD-tree anchoring over about 2,000 simulated S11 spectra, extended to 12\,\texttimes\,12-pixel layouts. \href{https://github.com/Galabavamsi/Antenna-Net}{GitHub}}

\project{Electron-GNN}{team, in progress, 2026}{Two-tower GATv2 graph network predicting absorption-spectrum peaks from molecular geometry; so far trained and checked on two small molecules (water and ammonia). \href{https://github.com/Galabavamsi/Electron-GNN}{GitHub}}

\section*{Experience}
\entry{Research Intern}{IIT Bhilai\sep Advisor: Dr. Arzad Alam Kherani\sep\href{https://github.com/Galabavamsi/RIM-SIM-V2}{RIS-SIM v2}\sep\href{https://github.com/Galabavamsi/RIS-SIM}{v1}}{May 2025 -- Jul 2025}
\begin{itemize}
  \item Built RIS-SIM, an open-source Python emulator for RIS-assisted radio links; designed RIS unit cells in CST Studio Suite (IEEE INDICON 2025, co-author). Rebuilt it in 2026 as RIS-SIM v2 (103 tests), then as sremu, an IQ-level emulator with about 2,100 tests (\href{https://github.com/Galabavamsi/sremu}{GitHub}).
\end{itemize}

\entry{Co-founder, Human Slop}{Anti-AI social app\sep\href{https://humanslop.in/}{Web app}\sep\href{https://play.google.com/store/apps/details?id=com.humanslop.app}{Play Store}}{2026 -- Present}
\begin{itemize}
  \item Co-building an Expo app (web and Android), a Fastify API, and PostgreSQL on Supabase; live on Google Play (50+ downloads). Added a keystroke-timing score and a device biometric check; both are heuristics, not proof of authorship.
\end{itemize}

\section*{Open Source}
\entry{Inspect Robots and inspect-robots-wandb}{Robocurve (YC S26)\sep\href{https://github.com/robocurve/inspect-robots}{Repository}\sep\href{https://pypi.org/project/inspect-robots-wandb/}{PyPI}}{Sep 2026 -- Oct 2026}
\begin{itemize}
  \item Evaluation framework for physical-AI policies. 2 merged PRs (a configurable LLM retry policy across 5 wire protocols; git state no longer logged as clean when unknown), 3 open. Wrote and maintain inspect-robots-wandb, a W\&B plugin on PyPI with a 100\% coverage gate and trusted publishing.
\end{itemize}

\entry{HFlow}{Hebbian Robotics (YC S26)\sep\href{https://github.com/Hebbian-Robotics/hflow}{Repository}}{Sep 2026 -- Oct 2026}
\begin{itemize}
  \item 8 merged PRs to an SDK that verifies robot training data: conversion keeps data whole or refuses it with a reason, e.g. depth videos that would silently drop from 12-bit to 10-bit (\href{https://github.com/Hebbian-Robotics/hflow/pull/400}{\#400}). A benchmark PR reported no safe speedup and changed nothing.
\end{itemize}

\section*{Publications}
\begin{itemize}
  \item \textbf{Co-author}~--- \textit{Inspect Robots: Evaluating the Capabilities and Safety of Embodied AI}, arXiv:2610.06306, 2026; accepted to the SPAIS Workshop at CoRL 2026 (poster). \href{https://arxiv.org/abs/2610.06306}{arXiv}
  \item \textbf{Co-author}~--- \href{https://ieeexplore.ieee.org/abstract/document/11392908/}{\textit{Experience with RF Energy Harvesting-Driven Self-Powered RIS}}, IEEE INDICON 2025.
  \item \textbf{Co-author}~--- \href{https://www.sciencedirect.com/science/article/pii/S1570870526001964}{\textit{A systematic literature review on simulation models and deployments for reconfigurable intelligent surfaces}}, Ad Hoc Networks (Elsevier), 2026.
  \item \textbf{Sole author (preprint draft, arXiv submission pending)}~--- \textit{Timed RIS Control in the Loop: A Sample-Accurate IQ-Level Emulator with Characterized Timing}, 2026. \href{https://github.com/Galabavamsi/sremu-paper/blob/master/preprint/sremu-preprint-draft.pdf}{Draft}
\end{itemize}

\section*{Skills}
{\raggedright\linepenalty=1000
\textbf{ML and data:} PyTorch, PyTorch Geometric, LightGBM, OpenCV, TFLite, NumPy, CuPy, Apache Arrow, Weights \& Biases, scikit-image, Qwen2-VL, ViT/RoBERTa, XLS-R/AASIST\par
\textbf{Languages:} Python, SQL, TypeScript/JavaScript, C/C++, Rust (open-source contributions), GLSL\par
\textbf{Tooling:} FastAPI, PostgreSQL, Docker, AWS, Git, Linux, GitHub Actions, pytest, Ruff, mypy, PyPI trusted publishing\par
\textbf{Scientific computing and simulation:} NS-3, MEEP (FDTD), CST Studio Suite, GNU Radio, MATLAB/Simulink, MuJoCo\par}

\section*{Achievements}
{\raggedright\linepenalty=1000
\textbf{5th of 232 teams}, WEAR Dataset Challenge @HASCA 2026 (Kaggle, final private leaderboard, macro-F1 0.93137)\sep
\textbf{AIR 200} (2026) and \textbf{AIR 700} (2025; 20,000+ participants), Amazon ML Challenge (team)\sep
\textbf{5th place}, Arista Networks Wi-Fi challenge, Inter IIT Tech Meet 14.0 (team,~Dec~2025)\sep
\textbf{Top 1\%}, JEE Advanced\sep
\textbf{PMSS} scholarship (2024)\sep
\textbf{Subroto Mukherjee} Scholarship (IAFBA, 2023)\sep
\textbf{YC Startup School}, Bengaluru (attended)\par}

\end{document}
